\section{Smoothing}
\begin{frame}{}
    \LARGE N-gram Models: \textbf{Smoothing}
\end{frame}

\begin{frame}{Zero Counts Are Not Zero Probabilities}
    \texttt{<UNK>} fixes words missing from $V$. It does nothing for a pair of words
    that are both in $V$ but never appeared \emph{together}.

    \vspace{1em}
    For the corpus \texttt{I am happy because I am learning} we computed
    \[
        P(\text{happy} \mid \text{I}) = \frac{C(\text{I happy})}{C(\text{I})}
        = \frac{0}{2} = 0
    \]
    even though ``I'' and ``happy'' are both perfectly ordinary words.

    \vspace{1em}
    \begin{itemize}
        \item One zero factor drives the whole sentence probability to $0$.
        \item $\log 0 = -\infty$, so perplexity is undefined rather than merely bad.
        \item Longer contexts make this worse: most trigrams a model meets at test
        time were never seen in training.
    \end{itemize}
\end{frame}

\begin{frame}[allowframebreaks]{Add-One (Laplace) Smoothing}
    \textbf{Idea:} pretend every possible n-gram was seen once more than it was.

    \[
        P_{\text{Laplace}}(w_i \mid w_{i-1})
        = \frac{C(w_{i-1}, w_i) + 1}{C(w_{i-1}) + |V|}
    \]

    \begin{itemize}
        \item The $+1$ in the numerator lifts unseen pairs off zero.
        \item The $+|V|$ in the denominator is what keeps each row summing to $1$:
        one unit is added for every word that could follow $w_{i-1}$.
        \item \textbf{Add-$k$} is the same move with a smaller nudge, $k < 1$:
        \[
            P_{\text{add-}k}(w_i \mid w_{i-1})
            = \frac{C(w_{i-1}, w_i) + k}{C(w_{i-1}) + k\,|V|}
        \]
    \end{itemize}

\framebreak
    \textbf{Worked example.} Same corpus, $V = \{$I, am, happy, because, learning$\}$,
    so $|V| = 5$ and $C(\text{I}) = 2$.

    \vspace{0.5em}
    \begin{columns}[t]
        \begin{column}{0.5\textwidth}
            \textbf{Unsmoothed}
            \[ P(\text{am} \mid \text{I}) = \frac{2}{2} = 1 \]
            \[ P(\text{happy} \mid \text{I}) = \frac{0}{2} = 0 \]
        \end{column}
        \begin{column}{0.5\textwidth}
            \textbf{Add-one}
            \[ P(\text{am} \mid \text{I}) = \frac{2+1}{2+5} = \frac{3}{7} \]
            \[ P(\text{happy} \mid \text{I}) = \frac{0+1}{2+5} = \frac{1}{7} \]
        \end{column}
    \end{columns}

    \vspace{0.5em}
    Nothing is zero any more, and the five smoothed probabilities still sum to $1$.
    The mass came out of the one bigram we actually observed, whose estimate fell
    from $1$ to $3/7$.

\framebreak
    \textbf{The trade-off:} that is a lot of mass to give away.

    \begin{itemize}
        \item Real vocabularies hold tens of thousands of words, so $|V|$ dwarfs
        $C(w_{i-1})$ and add-one hands almost all the probability to n-grams that
        were never seen.
        \item Add-$k$ softens this but leaves the same flaw: every unseen n-gram gets
        the \emph{same} estimate, whether it is plausible English or nonsense.
        \item Practical models instead fall back on shorter contexts, which are
        estimated from more data. If the trigram $(w_{i-2}, w_{i-1}, w_i)$ is unseen,
        use the bigram $(w_{i-1}, w_i)$; if that is unseen too, use the unigram.
        \begin{itemize}
            \item \textbf{Backoff:} drop to the shorter context only when the longer
            one has count zero.
            \item \textbf{Interpolation:} always mix all orders,
            $\hat{P} = \lambda_1 P(w_i) + \lambda_2 P(w_i \mid w_{i-1})
            + \lambda_3 P(w_i \mid w_{i-2}, w_{i-1})$, with
            $\sum_j \lambda_j = 1$ and the $\lambda_j$ tuned on held-out data.
        \end{itemize}
    \end{itemize}
\end{frame}
